% P17-D1: finite weak preference, ordinal representation and maximal choice.
\chapter{Finite Preferences and Ordinal Representation}
\label{sec:p17-finite-preferences}

A preference relation compares alternatives. An ordinal representation
assigns numbers that preserve those comparisons. This chapter develops the
finite case of Rubinstein's weak-preference and utility-representation
definitions \cite[2015 revision, pp.~5--8 and 12--14]{Rubinstein2015LectureNotes}.
It also constructs the induced choice correspondence on finite menus,
including ties \cite[p.~31]{Rubinstein2015LectureNotes}.

The statements concern a supplied relation on a fixed set of alternatives.
They do not establish that a person's observed comparisons satisfy the
assumptions, or that numerical differences measure intensity. All selected
results are proved in the annex. Empty sets, singletons and ties receive
explicit treatment.

\section{Weak and strict comparisons}
\label{sec:p17-preference-relations}

\begin{hbdefinition}{def:p17-weak-preference}{Finite weak preference}
Let \(X\) be a finite set of \emph{alternatives}. A binary relation
\(\succeq\subseteq X\times X\), read ``at least as preferred as'',
is a \emph{weak preference} when it satisfies:
\begin{enumerate}
\item \emph{Completeness:} for every \(x,y\in X\),
      \(x\succeq y\) or \(y\succeq x\), or both.
\item \emph{Transitivity:} for every \(x,y,z\in X\),
      \(x\succeq y\) and \(y\succeq z\) imply \(x\succeq z\).
\end{enumerate}
Completeness includes \(x=y\). A relation is \emph{reflexive} when
\(x\succeq x\) for every \(x\in X\); under the stated completeness
convention this follows, rather than being a separate assumption. The
convention is Rubinstein's Definition~2, with a different glyph for the
same weak comparison \cite[pp.~5--6]{Rubinstein2015LectureNotes}.

Define \emph{strict preference} and \emph{indifference} by
\[
\begin{aligned}
 (x\succ y)&\defeq
       (x\succeq y\text{ and not }y\succeq x),\\
 (x\sim y)&\defeq
       (x\succeq y\text{ and }y\succeq x).
\end{aligned}
\]
The notation \(\sim\) has this preference role throughout the chapter.
Indifference records two affirmative weak comparisons; it does not mean
that the alternatives have not been compared.
\index{preference!weak}\index{preference!strict}
\index{preference!indifference}
\end{hbdefinition}

\begin{hbdefinition}{def:p17-preference-quotient}{Indifference classes and their comparison}
For \(x\in X\), put
\[
 [x]\defeq\{y\in X:y\sim x\},\qquad
 \mathcal Q\defeq\{[x]:x\in X\}.
\]
On this set of classes define
\[
 (A\unrhd B)\defeq
   (\text{there exist }a\in A,\ b\in B\text{ with }a\succeq b).
\]
The next theorem shows that \(\mathcal Q\) is a partition of \(X\)
and that this comparison agrees with the comparison of any representatives.
We then write \(X/{\sim}\) for \(\mathcal Q\), the
\emph{indifference quotient}.

A \emph{total order} here is a reflexive, transitive, complete and
antisymmetric relation. Antisymmetry on classes means that
\(A\unrhd B\) and \(B\unrhd A\) imply \(A=B\).
It is not required of the original weak preference on alternatives.
\index{preference!indifference quotient}
\end{hbdefinition}

\begin{hbtheorem}{thm:p17-preference-structure}{proved=annex:p17.preference-structure}{The structure of finite weak preference}
A weak preference is reflexive. Its indifference relation is an equivalence
relation, and its classes partition \(X\). Strict preference is
irreflexive, asymmetric and transitive. For every \(x,y\in X\),
exactly one of
\[
 x\succ y,\qquad x\sim y,\qquad y\succ x
\]
holds. Moreover,
\[
 [x]\unrhd[y]\quad\Longleftrightarrow\quad x\succeq y,
\]
and \(\unrhd\) is a total order on \(\mathcal Q\). In particular,
comparisons between classes do not depend on the representatives chosen.
\end{hbtheorem}

\section{Finite ordinal representation}
\label{sec:p17-ordinal-representation}

\begin{hbdefinition}{def:p17-ordinal-representation}{Ordinal representation and class rank}
For a binary relation \(R\subseteq X\times X\), a function
\(u:X\to\mathbb{R}\) is an \emph{ordinal representation} when
\[
 xRy\quad\Longleftrightarrow\quad u(x)\geq u(y)
                         \qquad(x,y\in X).
\]
For weak preference, the source calls this a utility representation
\cite[p.~12]{Rubinstein2015LectureNotes}. The word ``ordinal'' emphasizes
that the displayed equivalence specifies comparisons, including ties.

For the indifference quotient of a weak preference, put
\[
\begin{aligned}
 k&\defeq|\mathcal Q|,\qquad I_k\defeq\{j\in\N:j<k\},\\
 r(A)&\defeq|\{B\in\mathcal Q:A\unrhd B\text{ and }B\ne A\}|,\\
 u_0(x)&\defeq r([x]).
\end{aligned}
\]
Thus \(r(A)\) counts classes strictly below \(A\), rather than
counting alternatives. We have \(I_0=\varnothing\); for \(k>0\),
\(I_k=\{0,\ldots,k-1\}\). Regard these natural numbers as real
numbers when using \(u_0\) as a representation.
\index{preference!ordinal representation}\index{utility!ordinal}
\end{hbdefinition}

\begin{hbtheorem}{thm:p17-finite-representation}{proved=annex:p17.finite-representation}{Finite representation and its degenerate cases}
A binary relation on a finite set has an ordinal representation exactly
when it is complete and transitive, with completeness as in
\cref{def:p17-weak-preference}.

For a weak preference, \(r:\mathcal Q\to I_k\) is a bijection and
\(u_0\) is an ordinal representation. Every representation \(u\)
also satisfies
\[
\begin{aligned}
 x\sim y&\quad\Longleftrightarrow\quad u(x)=u(y),\\
 x\succ y&\quad\Longleftrightarrow\quad u(x)>u(y).
\end{aligned}
\]
The empty set has one binary relation and one function to \(\mathbb{R}\),
both empty; that function represents that relation. A singleton has its
reflexive relation as its unique weak preference, represented by any real
value assigned to its sole alternative. If \(X\ne\varnothing\)
and all alternatives are indifferent, the representations are exactly the
constant functions, and \(u_0\) is constantly zero.
\end{hbtheorem}

\begin{hbmethod}{def:p17-class-elimination}{Constructing the ordered classes}
Given a weak preference on finite \(X\), form its indifference classes.
If no classes remain, finish. Otherwise select the least class in the
remaining total order, give its members the next value \(0,1,2,\ldots\),
remove that class, and repeat. A least class means one below every
remaining class. The proof of \cref{thm:p17-finite-representation} shows
that it exists and is unique whenever the remaining set is nonempty.
The procedure removes exactly \(k\leq|X|\) nonempty classes and
returns \(u_0\).

This is the class formulation of the source's finite elimination construction
\cite[pp.~13--14]{Rubinstein2015LectureNotes}. The source numbers its
successive groups from one; our ranks start from zero. For nonempty \(X\)
the values differ by one, as proved in the annex. No least or greatest
element of an empty set is asserted.
\end{hbmethod}

\section{Exactly what is unique}
\label{sec:p17-ordinal-uniqueness}

\begin{hbdefinition}{def:p17-increasing-scale-change}{Increasing changes of attained values}
For \(D\subseteq\mathbb{R}\), a function \(f:D\to\mathbb{R}\) is
\emph{strictly increasing} when \(s<t\) in \(D\) implies
\(f(s)<f(t)\). Write
\(u(X)\defeq\{u(x):x\in X\}\) for the attained values of a
function \(u:X\to\mathbb{R}\).
\end{hbdefinition}

\begin{hbtheorem}{thm:p17-ordinal-uniqueness}{proved=annex:p17.ordinal-uniqueness}{Uniqueness on the attained values}
Fix an ordinal representation \(u\) of a weak preference on finite
\(X\). A function \(v:X\to\mathbb{R}\) represents the same preference
exactly when there is a strictly increasing bijection
\[
 f:u(X)\longrightarrow v(X),\qquad v=f\circ u.
\]
When it exists, this \(f\) is unique. Because its domain is finite,
it extends to a strictly increasing function \(F:\mathbb{R}\to\mathbb{R}\).
Such an extension is not unique: its values outside \(u(X)\) are
not fixed by the representation.

In particular, every strictly increasing \(F:\mathbb{R}\to\mathbb{R}\) preserves
representation when composed with \(u\), as in the source's claim
\cite[p.~13]{Rubinstein2015LectureNotes}. The converse and the precise
finite-domain extension statement are proved here. The allowed changes
need not have the positive affine form \(t\mapsto\alpha t+\beta\),
where \(\alpha>0\).
\end{hbtheorem}

\section{Maximal choice from a finite menu}
\label{sec:p17-maximal-choice}

\begin{hbdefinition}{def:p17-maximal-choice}{Maximal choice correspondence}
For any binary relation \(R\) on \(X\) and any menu \(B\subseteq X\),
put
\[
 C_R(B)\defeq\{x\in B:xRy\text{ for every }y\in B\}.
\]
For a weak preference, \(C_{\succeq}\) is its \emph{maximal choice
correspondence}. An alternative \(x\in B\) is \emph{strictly
undominated in \(B\)} when there is no \(y\in B\) with \(y\succ x\).
The next theorem identifies these alternatives with \(C_{\succeq}(B)\).

On nonempty menus this is the source's induced choice correspondence
\cite[p.~31]{Rubinstein2015LectureNotes}; its required nonemptiness is
proved here from the finite weak-preference assumptions. We extend its
domain to the empty menu by the same formula, which gives
\(C_R(\varnothing)=\varnothing\). A correspondence can retain several
tied alternatives. Selecting one of them requires an additional choice rule.
\index{choice!maximal correspondence}
\end{hbdefinition}

\begin{hbtheorem}{thm:p17-maximal-choice}{proved=annex:p17.maximal-choice}{Maxima, representation and menu restriction}
Let \(\succeq\) be a weak preference on finite \(X\), and let
\(u\) represent it. For every \(B\subseteq X\),
\[
\begin{aligned}
 C_{\succeq}(B)
   &=\{x\in B:\text{no }y\in B\text{ satisfies }y\succ x\}\\
   &=\{x\in B:u(x)\geq u(y)\text{ for every }y\in B\}.
\end{aligned}
\]
If \(B\ne\varnothing\), this set is nonempty and is exactly the
intersection of \(B\) with the highest indifference class that \(B\)
meets. It is therefore unchanged by replacing \(u\) with another
representation. Every singleton menu \(\{x\}\) has choice set
\(\{x\}\), and the empty menu has empty choice set.

The correspondence recovers the comparisons by
\[
 x\succeq y\quad\Longleftrightarrow\quad
       x\in C_{\succeq}(\{x,y\}).
\]
Moreover, if \(B\subseteq A\subseteq X\) and
\(C_{\succeq}(A)\cap B\ne\varnothing\), then
\[
 C_{\succeq}(B)=C_{\succeq}(A)\cap B.
\]
\end{hbtheorem}

\section{Ties, cycles and the ordinal boundary}
\label{sec:p17-preference-examples}

\begin{hbdefinition}{def:p17-tied-preference-example}{A finite comparison with ties}
Let \(X=\{a,b,c,d,e\}\), with all five alternatives distinct, and
define \(t,v:X\to\mathbb{R}\) by the complete table
\[
\begin{array}{c|rrrrr}
 x&a&b&c&d&e\\\hline
 t(x)&0&0&1&2&2\\
 v(x)&0&0&1&10&10
\end{array}
\]
Define \(x\succeq_{\rm tie}y\) exactly when \(t(x)\geq t(y)\).
This table is an additional handbook example.
\end{hbdefinition}

\begin{hbtheorem}{thm:p17-tied-preference-example}{proved=annex:p17.tied-preference-example}{The classes, representations and choices in the tied example}
The relation \(\succeq_{\rm tie}\) is a weak preference. Its classes,
from lowest to highest, are
\[
 \{a,b\},\qquad\{c\},\qquad\{d,e\}.
\]
Its canonical representation \(u_0\) is \(t\), and \(v\) is a
second representation. Its choices include the following complete menu
calculations:
\[
\begin{aligned}
 C_{\succeq_{\rm tie}}(X)&=\{d,e\},\\
 C_{\succeq_{\rm tie}}(\{a,c,d,e\})&=\{d,e\},\\
 C_{\succeq_{\rm tie}}(\{a,b,c\})&=\{c\},\\
 C_{\succeq_{\rm tie}}(\{a,b\})&=\{a,b\}.
\end{aligned}
\]
In addition every singleton chooses its sole member, and the empty menu
has empty choice set.
\end{hbtheorem}

\begin{hbdefinition}{def:p17-preference-cycle}{A complete cyclic comparison}
On \(X_{\rm cyc}=\{a,b,c\}\), with the alternatives distinct, define
\[
 R_{\rm cyc}\defeq
 \{(a,a),(b,b),(c,c),(a,b),(b,c),(c,a)\}.
\]
These are all its ordered pairs. This is a comparison relation to be
checked against the preference assumptions.
\end{hbdefinition}

\begin{hbtheorem}{thm:p17-preference-cycle}{proved=annex:p17.preference-cycle}{A cycle prevents ordinal representation}
The relation \(R_{\rm cyc}\) is reflexive and complete but is not
transitive. It has no ordinal representation, and
\[
 C_{R_{\rm cyc}}(X_{\rm cyc})=\varnothing.
\]
Thus completeness alone does not give the representation or nonempty
maximal choice guaranteed for weak preferences.
\end{hbtheorem}

\begin{hbtheorem}{thm:p17-ordinal-boundary}{proved=annex:p17.ordinal-boundary}{The same ordering permits different gaps and averages}
For the two representations in the tied example there are no real numbers
\(\alpha>0,\beta\) with \(v(x)=\alpha t(x)+\beta\) for all
\(x\in X\). Their adjacent differences and arithmetic averages satisfy
\[
\begin{aligned}
 t(c)-t(a)&=t(d)-t(c)=1,\\
 v(c)-v(a)&=1,\qquad v(d)-v(c)=9,\\
 \frac{t(a)+t(d)}2&=t(c),\\
 \frac{v(a)+v(d)}2&=5>v(c).
\end{aligned}
\]
Both representations nevertheless give exactly the same comparisons and
maximal choices on the original alternatives.
\end{hbtheorem}

\begin{hbwarning}{def:p17-ordinal-interpretation-boundary}{Ordinal values and additional interpretations}
An ordinal representation preserves ordering and indifference. The represented
ordering alone does not fix gaps between numerical values or comparisons
formed by averaging them. Consequently an intensity scale or an
expected-utility interpretation does not follow from this finite ordering.
No lottery domain, probabilities or preferences over lotteries have been
introduced here. Such a representation requires its own assumptions and
proofs.

The source's countable and infinite-domain results, continuity conditions,
choice-rationalization theorems and expected-utility treatment are separate
from this chapter's finite foundation. This unit also does not establish
empirical adequacy of completeness or transitivity, or complete the broader
P17 preference and decision-theory subject.
\end{hbwarning}
